\documentclass[11pt,reqno]{amsart}

\usepackage[margin=1in]{geometry}
\usepackage[T1]{fontenc}
\usepackage{lmodern}
\usepackage{microtype}
\usepackage{mathtools}
\usepackage{amssymb}
\usepackage{braket}
\usepackage{xcolor}
\usepackage[colorlinks=true,linkcolor=blue!55!black,citecolor=blue!55!black,
  urlcolor=blue!55!black]{hyperref}

\theoremstyle{plain}
\newtheorem{theorem}{Theorem}[section]
\newtheorem{corollary}[theorem]{Corollary}
\newtheorem{proposition}[theorem]{Proposition}
\newtheorem{lemma}[theorem]{Lemma}
\newtheorem{conjecture}[theorem]{Conjecture}
\theoremstyle{remark}
\newtheorem{remark}[theorem]{Remark}

\usepackage[nameinlink,capitalize,noabbrev]{cleveref}

\DeclareMathOperator{\Tr}{Tr}
\DeclareMathOperator{\rank}{rank}
\newcommand{\cL}{\mathcal L}
\newcommand{\ketbra}[2]{\lvert #1\rangle\!\langle #2\rvert}
\newcommand{\ip}[2]{\langle #1,#2\rangle}
\newcommand{\norm}[1]{\lVert #1\rVert}
\newcommand{\abs}[1]{\lvert #1\rvert}
\newcommand{\C}{\mathbb C}

\title{The higher Choi constants of the double-skew space}
\author{Elazar Gershuni}
\date{\today}

\begin{document}
\begin{abstract}
For the double-skew space $\mathfrak{so}(U)\otimes\mathfrak{so}(V)$ we
determine the quantities
$\kappa_k=\max\{\sum_{j\le k}s_j(K)^2:\norm K_F=1\}$, equivalently the
Johnston--Kribs $S(k)$-norms of the orthogonal projection onto it, equivalently
the Ky Fan thresholds $a_k$ of M\l ynik--Osaka--Marciniak for that space.  We
show $\kappa_k=\min\{k,4\}/4$ for every $k\ge2$, so that the associated Choi
constant is $\beta_k=\min\{1,4/k\}$; and we show that this formula
\emph{fails} at $k=1$, by exhibiting an explicit element attaining
$\kappa_1\ge(m-1)/2m$ where $m=\min\{\dim U,\dim V\}$.  Numerically that value
is exact.
\end{abstract}
\maketitle

\section{Setting}

Let $U$ and $V$ be finite-dimensional with
$\dim U,\dim V\ge2$.  Fix orthonormal bases and write
$\mathfrak{so}(W)=\{L:L^T=-L\}$, and let
\[
    \mathcal A:=\mathfrak{so}(U)\otimes\mathfrak{so}(V)\subseteq\cL(U\otimes V)
\]
be the linear span of the elementary tensors $L\otimes M$.  For $k\ge1$ put
\begin{equation}
\label{eq:kappa}
    \kappa_k
    :=
    \max\Bigl\{\textstyle\sum_{j\le k}s_j(K)^2
      :K\in\mathcal A,\ \norm K_F=1\Bigr\},
\end{equation}
$s_1\ge s_2\ge\cdots$ being singular values.  Three readings coincide.  By the
$k$-th case of Johnston--Kribs projection duality, $\kappa_k$ is the $S(k)$-norm
of the Frobenius-orthogonal projection $\Pi$ onto $\mathcal A$.  By
M\l ynik--Osaka--Marciniak, $4\kappa_k$ is the threshold $a_k$ characterizing
$k$-positivity of $a\Tr(\cdot)I-\Phi_{\mathcal A}$.  And in the normalization in
which the Kraus operators $L_\alpha\otimes M_\beta$ of
$\Lambda_U\otimes\Lambda_V$ have $\norm{\cdot}_F^2=4$, the Choi constant is
\begin{equation}
\label{eq:beta-kappa}
    \beta_k(\Lambda_U\otimes\Lambda_V)=\frac{4\kappa_k}{k}.
\end{equation}

The single input below is the pair bound established in the proof of the
double-skew action lemma, in the form proved there --- for \emph{every} pair,
not merely for the two largest:
\begin{equation}
\label{eq:pair}
    s_a(K)^2+s_b(K)^2\le\tfrac12\norm K_F^2
    \qquad\text{for all }a\ne b,\ K\in\mathcal A .
\end{equation}

\section{\texorpdfstring{$k\ge2$}{k >= 2}}

\begin{theorem}
\label{thm:kappa}
For every $k\ge2$, $\kappa_k=\min\{k,4\}/4$, and hence
$\beta_k(\Lambda_U\otimes\Lambda_V)=\min\{1,4/k\}$.
\end{theorem}

\begin{proof}
Normalize $\norm K_F=1$.

\emph{Upper bounds.}  Summing \cref{eq:pair} over all $\binom k2$ pairs drawn
from $\{1,\ldots,k\}$, and using that each index occurs in exactly $k-1$ of
them,
\[
    (k-1)\sum_{j\le k}s_j^2
    =\sum_{a<b\le k}\bigl(s_a^2+s_b^2\bigr)
    \le\binom k2\cdot\tfrac12
    =\frac{k(k-1)}4 ,
\]
so $\sum_{j\le k}s_j^2\le k/4$ for every $k\ge2$.  Together with the trivial
$\sum_{j\le k}s_j^2\le1$ this gives $\kappa_k\le\min\{k,4\}/4$.

\emph{Attainment.}  A single elementary tensor $L_{ab}\otimes M_{cd}$, with
$L_{ab}=E_{ab}-E_{ba}$, has exactly four nonzero singular values, all equal, so
after normalization $\sum_{j\le k}s_j^2=\min\{k,4\}/4$ for every $k$.
\end{proof}

The assumptions $\dim U,\dim V\ge2$ are necessary for the statement as
written.  If either factor has dimension below two, then
$\mathcal A=\{0\}$, so the Frobenius-unit set in \cref{eq:kappa} is empty and
the displayed maximum is not defined; the associated completely positive map
is zero.

\begin{remark}
\label{rem:everypair}
The averaging step is where \cref{eq:pair} must hold for \emph{every} pair.  The
weaker statement for the two largest singular values gives, by grouping into
consecutive pairs, only $\sum_{j\le k}s_j^2\le\lceil k/2\rceil/2$ --- that is
$k/4$ for even $k$ but $(k+1)/4$ for odd $k$, hence the trivial bound at $k=3$.
The every-pair form was proved for an unrelated reason, namely to avoid ordering
the singular values before the last step of the double-skew lemma.
\end{remark}

\begin{remark}
\label{rem:action}
Nothing above needs singular values.  Since $\sum_{j\le k}s_j(K)^2$ is the
maximum of $\sum_{i\le k}\norm{Kx_i}^2$ over orthonormal $k$-tuples, and
\cref{eq:pair} is exactly the two-vector action bound, the same averaging gives
directly
\[
    \sum_{i=1}^k\norm{Kx_i}^2\le\frac k4\norm K_F^2
    \qquad\text{for every orthonormal }x_1,\ldots,x_k ,
\]
which is the form the amplification argument consumes and the form in which the
statement is machine-checked.
\end{remark}

\section{\texorpdfstring{$k=1$}{k = 1} is different}

Extrapolating \cref{thm:kappa} to $k=1$ would give $\kappa_1=\frac14$ and
$\beta_1=1$.  Both are false as soon as both factors have dimension at least
three.  Write $m:=\min\{\dim U,\dim V\}$.

\begin{theorem}
\label{thm:k1}
Let $\dim U=\dim V=m\ge2$, let $\{L_\alpha\}$ be a Frobenius-orthonormal basis
of $\mathfrak{so}(m)$, and set
\[
    K_m:=\sum_\alpha L_\alpha\otimes L_\alpha\in\mathcal A .
\]
Then
\begin{equation}
\label{eq:Km}
    K_m=\tfrac12\bigl(m\ketbra\omega\omega-F\bigr),
\end{equation}
where $F$ is the flip on $\C^m\otimes\C^m$ and
$\omega=m^{-1/2}\sum_ae_a\otimes e_a$.  Consequently $K_m$ has eigenvalue
$\frac{m-1}2$ with multiplicity one and eigenvalues $\pm\frac12$ on the
orthogonal complement, so
\[
    \frac{s_1(K_m)^2}{\norm{K_m}_F^2}
    =\frac{(m-1)^2/4}{m(m-1)/2}
    =\frac{m-1}{2m},
    \qquad\text{whence}\qquad
    \kappa_1\ge\frac{m-1}{2m} .
\]
In particular $\kappa_1>\frac14$ and $\beta_1=4\kappa_1>1$ for every $m\ge3$.
\end{theorem}

\begin{proof}
Take $L_\alpha=(E_{ab}-E_{ba})/\sqrt2$ for $a<b$.  The summand of
$\sum_{a<b}(E_{ab}-E_{ba})\otimes(E_{ab}-E_{ba})$ is invariant under $a\leftrightarrow b$,
so
\[
    K_m
    =\tfrac12\sum_{a<b}(E_{ab}-E_{ba})^{\otimes2}
    =\tfrac14\sum_{a,b}(E_{ab}-E_{ba})^{\otimes2}
    =\tfrac12\Bigl(\sum_{a,b}E_{ab}\otimes E_{ab}
                  -\sum_{a,b}E_{ab}\otimes E_{ba}\Bigr).
\]
The second sum is $F$.  For the first, $E_{ab}e_c=\delta_{bc}e_a$ gives
$\sum_{a,b}E_{ab}\otimes E_{ab}\,(e_c\otimes e_d)=\delta_{cd}\sum_ae_a\otimes e_a$,
which is $m\ketbra\omega\omega$.  This is \cref{eq:Km}.

Now $F$ is $+1$ on $\operatorname{Sym}^2$ and $-1$ on $\wedge^2$, and
$\omega\in\operatorname{Sym}^2$ with $F\omega=\omega$.  Hence $K_m$ acts as
$\frac12(m-1)$ on $\C\omega$, as $-\frac12$ on
$\operatorname{Sym}^2\ominus\C\omega$, and as $+\frac12$ on $\wedge^2$.  Its
squared Frobenius norm is therefore
\[
    \frac{(m-1)^2}4+\frac{m^2-1}4=\frac{m(m-1)}2,
\]
and $s_1=\frac{m-1}2$ since $\frac{m-1}2\ge\frac12$ for $m\ge2$.  Dividing gives
the stated ratio.  For $m\ge3$ it exceeds $\frac14$; unequal factor dimensions
only enlarge $\mathcal A$, so the bound persists with $m=\min\{\dim U,\dim V\}$.
\end{proof}

\begin{conjecture}
\label{conj:k1}
$\kappa_1=(m-1)/2m$, so that $\beta_1=2(m-1)/m$.
\end{conjecture}

Alternating maximization of $\norm{\Pi(M)}_F^2$ over rank-one $M$ agrees to five
decimals at
$(\dim U,\dim V)\in\{(2,2),(3,2),(5,2),(3,3),(4,3),(5,3),(4,4),(5,5),(6,6)\}$,
in each case reproducing $(m-1)/2m$ with $m=\min\{\dim U,\dim V\}$.  Only the
upper bound is open: \cref{eq:pair} applied to a pair $(1,b)$ with $s_b=0$ gives
$\kappa_1\le\frac12$, and $(m-1)/2m\to\frac12$, so that estimate is
asymptotically sharp.

\begin{remark}[Why the extremizer changes]
An elementary tensor $L\otimes M$ has four \emph{equal} singular values, which
is what makes it optimal for every $k\ge2$ at once.  At $k=1$ the mass should
instead be concentrated in a single singular value, and $K_m$ does that: it
places $(m-1)^2/4$ of the total $m(m-1)/2$ on one eigenvector, the invariant
$\omega$.  Equivalently, the top singular value of an element of $\mathcal A$ is
not forced to be doubly degenerate --- an assumption which, if it held, would
turn \cref{eq:pair} into $\kappa_1\le\frac14$.  It fails from $m=3$.
\end{remark}

\begin{remark}[On product testers]
For $u=a\otimes b$ and $v=c\otimes d$ one has
$\Pi(\ketbra uv)=P_\wedge(\ketbra ac)\otimes P_\wedge(\ketbra bd)$ with
$P_\wedge(X)=\frac12(X-X^T)$, and the resulting ratio never exceeds $\frac14$.
So the extremizers of \cref{thm:k1} are necessarily entangled across the
$U:V$ cut, and $m=2$ is exactly the case where product testers suffice.
\end{remark}

\end{document}
